\subsection{Wurzeln}
\label{subsec:fundamentals_roots}

\begin{samepage}
  \begin{DefinitionBox}{$n$-te Wurzel}{fundamentals_nth_root}
    Wurzelziehen kehrt das Potenzieren um: Bei der $n$-ten Wurzel aus $a$ wird
    der Wert $b$ gesucht, dessen $n$-te Potenz wieder $a$ ergibt.\par\smallskip
    \noindent
    \begin{minipage}[c]{0.54\linewidth}
      \[
        \sqrt[n]{a}=b
        \quad \Longleftrightarrow \quad
        b^n=a,
        \qquad n\in\N.
      \]
      Für $n=2$ schreiben wir kurz: $\sqrt a := \sqrt[2]{a}$.
    \end{minipage}\hfill
    \begin{minipage}[c]{0.4\linewidth}
      \centering
      \CachedTikzFigure{content/resources/chapters/fundamentals/tikz/root_anatomy.tex}
    \end{minipage}\par
  \end{DefinitionBox}

  Damit gilt insbesondere \[ \sqrt[n]{a^n}=a \qquad n\in\N. \]

  \vspace{-0.5\baselineskip}
  \begin{ExampleBox}{Wurzeln bestimmen}{fundamentals_even_odd_roots}
    \begin{enumerate}[label=(\alph*),itemsep=0.55em]
      \item $\sqrt{9}\Cloze{2cm}{=3}$,
        denn \Cloze{6cm}{$3^2=9$}.
      \item $\sqrt[3]{-8}\Cloze{2cm}{=-2}$,
        denn \Cloze{6cm}{$(-2)^3=-8$}.
      \item $\sqrt{-9}\Cloze{2.0cm}{\notin\R}$,
        denn \Cloze{6cm}{$b^2\geq0$ für alle $b\in\R$}.
    \end{enumerate}
  \end{ExampleBox}
\end{samepage}

\begin{KeyPointBox}{Reelle Wurzeln}{fundamentals_real_roots}
  Bei negativen Radikanden $a<0$ entscheidet der Wurzelexponent:
  \[
    \begin{array}{c@{\qquad}l}
      n\text{ ungerade}: &
        \Cloze{9cm}{\sqrt[n]{a} \in \R \quad \text{und} \quad \sqrt[n]{a} < 0.}\\[0.35em]
      n\text{ gerade}:   &
        \Cloze{9cm}{\sqrt[n]{a} \notin \R .}
    \end{array}
  \]
  Bei geradem Wurzelexponenten bezeichnet das Wurzelzeichen stets den
  nichtnegativen Wert, zum Beispiel
  \[
    \Cloze{8cm}{\sqrt{9}=3
      \qquad \text{aber } \sqrt{9} \neq \pm3.}
  \]
\end{KeyPointBox}
